% 原书第 259 页（PDF 第 272 页）
\BookPage{259}{272}
\BookHeading{论清代回疆边外浩罕强权崛起的中国因素}{论清代回疆边外浩罕强权崛起的中国因素}{陈旺城*\\[2pt]段昌国**}
\BookSubheading{一、前言}

「浩罕」（Khokand)，清朝文献上另译作「霍罕」，「豪罕」或「敖罕」1。这一伊斯兰强权是十八世纪至十九世纪才开始强大的一个以中亚费尔干纳（Farghana)盆地为中心，由鸟兹别克明格（UzbegMing）部所建立的政权。根据早期中亚的穆斯林史料，如《昔班尼传》(Shaybani-name)，其中曾谈及明格部的情形。而《阿布海尔史》(Tarikhi Abukhayr Khani)2中说:「金帐、白帐两汗国解体后，『游牧月即别』（作者案：即Uzbeg之音译，另译作乌兹别克）诸部分属右翼、左翼和中翼。其中名称为Tümän一Ming人者，为中翼的五部之一。」3而有关「Tümän」一词，潘志平氏引《哈萨克诸汗国史料集》一书，页494，认为「是个具有一定意义的社会术语，其意为『万人』，『万人军队』，『可提供万人军队的地区』」4。如此解释，应该是可信的，因为自匈奴头曼（Tümän）以来，一直到现在突回族人、蒙古人等，依然是这个意思。「Ming」是突回语「千」的意思，「Tümän一Ming」是「万千」的意思。内涵为「拥有万千之军队者」。所以说，十五世纪明格部是术赤五子昔班的六世孙阿布海尔汗(AbukhayrKhan)的月即别国家中的重要部族之一。若要更往前推，明格部可以说是出自术赤的钦察系，原活动于里海，后为「游牧月即别人」一支，

\vfill
\begin{BookNotes}
\BookNote{*}{作者任职于佛光大学生命与宗教学系。}
\BookNote{**}{作者任职于佛光大学未来学系。}
\BookNote{1.}{傅恒，《钦定皇舆西域图志》卷45（台北，文海出版社，民国59年)，页12。}
\BookNote{2.}{卡马丁·比纳（1453-1512)，《昔班尼传》，书成于1505-1507之间；马斯乌德·乌斯曼，《阿布海尔史》，书成于1543-1544年间。详见艾哈迈多夫，《中亚历史地理文献》(塔什干，1895年)，页20、54、77、107，引自潘志平，〈中亚浩罕明格朝前史初探〉，《西北史地》三期（兰州，兰州大学出版社，1987年），页62。}
\BookNote{3.}{艾哈迈多，《游牧月即别国家》（莫斯科，1965年)，页106，引自潘志平，前揭文，页62。}
\BookNote{4.}{同上注，潘志平，前揭文。}
\end{BookNotes}

% 原书第 260 页（PDF 第 273 页）
\BookPage{260}{273}
\noindent 「十六世纪随阿布海尔汗的昔班尼（作者案：为阿布海尔汗之孙）来到河中地区而遍及中亚的」5。

此一由明格部人所建立的国家，后发展成为以费尔干纳为中心，最强盛的时候，其疆域几乎占据了整个锡尔河（Syr-Darya）流域，是当时回疆边外中亚最大且最重要的伊斯兰强权之一。现在为了说明这一曾经是清代回疆变作事件主要支持者，且与清代中国发生过密切政治与经济关系的浩罕伊斯兰强权崛起情形，分析其中国背景因素如下。

\BookSubheading{二、清朝平定准部，浩罕枷锁去除}

清代所称之「卫拉特」或「厄鲁特」均为原蒙元时期，义为林木中百姓的「斡亦刺」或明代时所称之「瓦刺」，均为历史上不同时期「Oirat」一词的音转和异译。而准噶尔（Dzungar)之称呼，为蒙古语「左翼」的意思。因为明末清初，准噶尔部乃居新疆西蒙古四部（准噶尔、杜尔伯特、和硕特、土尔扈特）之左而得名。其后在十七世纪四、五十年代，始由准噶尔部的首领巴图鲁浑台吉（1643-1653）的领导下统一而强大起来，所以「准噶尔部」一词，才由一个游牧民族名称转变成为清代厄鲁特各部的统称。当时该部势力「并控制了塔什干（Tashkend）、安集延、撒马尔干等贸易中心」6，及至巴图鲁浑台吉六子噶尔丹（1644-1679）统治时，进行扩张兼并政策，也在「1683、1684、1685年还征讨了吉尔吉斯人和费尔干纳人」7，且把准噶尔的政治中心转移到伊犁河谷，除辖制天山南部外，并「将自己的势力扩展到撒马尔罕、布哈拉、乌尔根齐地区」8。

继噶尔丹之后，其侄策妄阿布坦（1697-1727）及策妄阿布坦子噶尔丹策零（1727-1745）继立为「台吉」。因为注重生产发展，所以较噶尔丹时期，具有着某

\vfill
\begin{BookNotes}
\BookNote{5.}{同上注。}
\BookNote{6.}{中国社会科学院民族研究所编，《准噶尔史略》（北京，人民出版社，1985年），页65。}
\BookNote{7.}{潘志平，《中亚浩罕国与清代新疆》（北京，中国社会科学出版社，1991年)，页91。}
\BookNote{8.}{同上注，页93。}
\BookNote{9.}{台吉（Taidji)一词为蒙古贵族爵位之名称，为汉语「太子」之讹转。详见林恩显，《中国边政通论》(台北，华泰书局，民国78年)，页212。而刘义棠则以为是突回语Tegin之口语音译，以往译作「特勤」、「的斤」，乃可汗子弟、贵族爵位名称，详见刘义棠中国边疆民族史，上课笔记。}
\end{BookNotes}

% 原书第 261 页（PDF 第 274 页）
\BookPage{261}{274}
\noindent 种程度的相对稳定与和平，但也仍经常与其周遭部族发生冲突与争战。而当时费尔干纳沙鲁赫所创立的浩罕明格政权，成立初期，仍极微弱，因此也经常处于准噶尔势力的威胁之下。康熙五十三年（1714），准噶尔部入侵过费尔于纳。依《大清历朝实录·圣祖实录》所载，策妄阿布坦的儿子就曾率兵攻打安箭（即安集延），「其子往攻安箭地方之布鲁特，被杀者五百人，回时又多染疾而死」10。又「1723年，准噶尔的军队入侵塔拉斯、塔什干地区，主要是同哈萨克人交战，但他的军队也一度进入费尔干纳」11，十八世纪三十年代初，「准噶尔的军队再次深入到天山山区以南，进入帕米尔，甚至深入到巴达克山，强行带走了巴达克山的贵族作为人质」12。此举，费尔干纳的明格政权显然也受到骚扰。十八世纪的四十年代中期，准部之入侵费尔干纳，影响更为深远。由可能作为职业军人，且亲身参加了（1826）张格尔（Jhanger)入侵喀什噶尔（Kashgar）活动，自称是受浩罕末代汗-胡达雅尔（Khudayar，1845-1858,1862,1865-1875）之命而作,成书在1871-1872年间,主要在讲述1721-1872年间的浩罕历史的毛拉·尼亚孜·穆罕默德的《沙鲁赫史》(Tarikhi Shahrukhi）中即记载：「巴巴伯克是热依姆汗（即阿卜都尔热依姆）的长子，当阿卜都尔噶里木称速檀时，卡尔梅克的军队从卡伊纳尔(Kaynar)开来，包围了浩罕城堡。由于没有力量抵抗，阿卜都尔噶里木被迫把巴巴伯克作为人质交出来。」13根据俄文资料，「1745-1746年准噶尔人同阿卜都尔噶里木交战」14，此事清朝史料，记述尤为详实，「驻藏副都统傅清奏。遵旨饬令郡王颇罗鼐，于今年卡隘处，派素所深信之人，远出侦察准噶尔踪迹。十月初三日，据颇罗鼐呈称，本年三月二十日派深信之萨嘉达克齐、搏尔搏等三十名，探望准噶尔做善事回巢踪迹…又有阿木宾禅汗，强暴有力，占据四处部落，所属有阿卜都尔噶里木部落。中途又有博洛尔部落。博洛尔部落内无汗，系哈屯楚哈尔办事。噶尔丹策零在阿卜都尔噶里木看视行兵之路，遣回子巴克达呼雅逊都喇呼沙、巴勒宾等四十人前往。彼处人猜疑，杀死三十三人，惟七人遁归。噶尔丹策零愈怒，派厄鲁特兵二万四千，哈萨克兵四千、吉尔吉斯兵二千，令

\vfill
\begin{BookNotes}
\BookNote{10.}{《清圣祖实录》卷263（台北，华文书局，民国53年），页23下-25上，康熙五十四年五月壬子，议政大臣等奏。}
\BookNote{11.}{伊凡诺夫，《中亚史纲》（莫斯科，1958年)，页111，引自潘志平，前揭书，页41。}
\BookNote{12.}{潘志平，前揭书，页41。}
\BookNote{13.}{同上注。}
\BookNote{14.}{同上注。}
\end{BookNotes}

% 原书第 262 页（PDF 第 275 页）
\BookPage{262}{275}
\noindent 伊族台吉卜登图拉尔、满积图拉尔、衮都尔图拉尔，于今年正月内起程前往。但闻阿卜都尔噶里木汗力量充足，兵丁强健，胜负未知如何？」15

在阿卜都尔噶里木死后，浩罕宫廷发生权利斗争时，准噶尔部也试图干预，《沙鲁赫史》中记载：「在得知阿卜都尔噶里木死后，卡尔梅克人把巴巴伯克派回浩罕说:『要知道他是热依姆汗的长子，他应得到政权。』浩罕接受了这一指令，让他登王位将近一年。不到一年，高级官吏将军对聚集到别什、阿雷克（Bish一Arek），将巴巴伯克杀死在那里，额尔德尼重新称汗。」16巴巴伯克的被杀，显现准噶尔的没落。

清朝原本对于这个十八世纪初期由浩罕沙鲁赫伯克所创建的明格政权，被当时正称雄中亚的准噶尔部侵凌，所持的只是「本朝并不之问」的态度。直到乾隆十年(1745)至十一年(1746)准部对阿卜都尔噶里木发动战争，才引起清朝的关注。而布哈拉到十八世纪上半叶也开始衰弱，致使其国内许多乌兹别克部族的「伯克」18在各地已自行其事。浩罕所处的费尔干纳盆地与准部为近邻，名义上虽仍为布哈拉汗国的一部分，但却不时的遭受到准噶尔的侵侮。甚至一度曾为准噶尔所降伏而成为其所属。直至乾隆二十四年（1759）清朝平定准、回两部，统一新疆后，才使浩罕得以完全摆脱对布哈拉的依附，并解除长期以来，屡受准噶尔部所加诸于其部族身上的侵扰与桎梏，而走上真正独立的道路。

\vfill
\begin{BookNotes}
\BookNote{15.}{《清高宗实录》卷252，页18下-20下，乾隆十年十一月乙亥条。}
\BookNote{16.}{潘志平，前揭书，页42。}
\BookNote{17.}{《清圣祖实录》卷183，页4，康熙三十六年四月甲寅谕。}
\BookNote{18.}{伯克（Beg）一词，为突回语音译，其意义为贵族、官吏等特权阶级的代表，官吏的总称。其成为一种制度是在清乾隆年间。功能在「为清廷维持地方的安定与征收赋税，派遣徭役。另一方面在维持回部原有传统社会文化」等。详见刘义棠，《维吾尔研究》（台北，中正书局，民国64年），页282-283；及林恩显，《清朝在新疆的汉回隔离政策》（台北，商务印书馆，民国77年），页93。}
\end{BookNotes}

% 原书第 263 页（PDF 第 276 页）
\BookPage{263}{276}
\BookSubheading{三、宗主清朝，浩罕伊斯兰政权刍形建立}

在十七一十八世纪的时候，因为准噶尔部的强大，割据了中国西北边疆的一隅，加之战乱纷繁，交通阻隔，致使清代王朝在开国一百多年之后，才与浩罕所处的葱岭以西的费尔干纳地区开始进行直接的官方关系。乾隆二十三年（1758)，清平回疆大小和卓之役，初期并不顺利，苦攻库车数月而成，后来兆惠先锋部队也在叶尔羌受困黒水营。为了使战事顺利进行的同时，清政府亦在中亚展开外交活动。当时参赞大臣富德遣侍卫蒙固尔岱等赴塔什罕地区宣谕威德。其时，在塔什罕的莫尔多萨木什、沙达和吐尔占三「和卓」19，均羁属于哈萨克右部。后来因为「吐尔占排斥哈萨克所置和卓莫尔多萨木什，哈萨克因此兴兵问罪。蒙固尔岱谕以睦邻守土之道，平息了他们的争端」20，于是哈萨克右部图里拜、塔什罕和卓吐尔占、莫尔多萨木什等向清表示归诚，「哈萨克、布鲁特，系属外藩，伊等入觐，虽藉抒诚悃，亦为希冀赏赐，如果情词恳切，自应准其来京，候朕酌量加恩」21。侍卫乌尔登、托伦泰亦受命往抚布鲁特等，且「招降安集延、纳木干二城回人」22。直至乾隆二十四年（1759）夏，回众节节败阵，为追击大小和卓，截堵其逃路，兆惠将军因此派遣侍卫达克塔纳「传檄额尔德尼伯克，晓以利害」23。额尔德尼迎接清使入城，盛情款待，「至其境，额尔德尼酋，迎至城内，日馈羊、酒、瓜果、糇粮、良马，诣访中国疆域、物产、风俗形势、兵马器械、侍卫广宣朝廷威德，额尔德尼畏慕、奉表」24请内附，「额尔德尼伯克等请大皇帝钦差统兵、至威至勇将军等安。仰赖上天之德、大皇帝洪福，我等群生亦获安业。闻将军至叶尔羌、喀什噶尔后，布鲁特、哈萨克等众部落俱传檄安定，而我部檄书未至，日夜以冀。幸侍卫赉将军书来，我等得为大皇帝臣仆，永受康宁，不

\vfill
\begin{BookNotes}
\BookNote{19.}{和卓一词，源于波斯语khoja，国际音标Xodja，汉文译音中有和加、和扎、霍加、火者、噶查…。在中亚一带有「高贵者」，「富有的商人」的含意。羽田明氏则指原为萨曼王朝的一重要官职名称，后成为贵族的称谓。但在中国回疆地区，则指神秘教派（Nakshbandiye)的圣裔，具有与「赛伊德(Sayyid)一穆罕默德之女蒂玛与阿里的后裔等同的含意。}
\BookNote{20.}{潘志平，前揭书，页45。}
\BookNote{21.}{《清高宗实录》，卷571，页23，乾隆二十三年九月辛亥富德奏。}
\BookNote{22.}{同上注，卷575，页38，乾隆二十三年十一月辛亥条。}
\BookNote{23.}{同上注，卷588，页10，乾隆二十四年六月乙卯条。}
\BookNote{24.}{魏源，《圣武记》卷4(台北，中华书局，民国51年)，页27。}
\end{BookNotes}

% 原书第 264 页（PDF 第 277 页）
\BookPage{264}{277}
\noindent 胜感悦」25。从此浩罕与清朝即展开了广泛的政治、商贸接触，并开启了清朝与浩罕的宗主外藩关系。

乾隆二十四年（1759）底，浩罕即派出了第一个朝觐团进京，这是浩罕的第一次觐使团，清政府对此相当重视，高宗下令「一路台站、马匹自应整理妥协」26，以迎送使团。根据《清实录》记载，乾隆二十四年至嘉庆十五年（1810），浩罕三代伯克（额尔德尼、纳尔巴图、爱里木）九次派出使节进京朝觐清朝皇帝（乾隆二十四年、二十七年、三十二年、三十五年、四十七年、五十七年、嘉庆元年、七年）直至嘉庆十四年（1809）的这一次例行朝觐，因清廷发现，所译出爱里木伯克呈文中有「语欠恭顺」，「殊失属臣之理」而传谕来使不令进京。后经当时喀什噶尔参赞大臣那彦成复查旧案，检查了自乾隆二十四年（1975）浩罕通使至今（嘉庆十四年）共二十三次，（即乾隆二十四年、二十六年、三十年二次、三十一年、三十二年、三十四年三次、三十五年、三十六年、三十七年、三十八年、四十一年、四十七年、四十八年、五十三年、五十六年、五十七年、六十年、嘉庆元年、四年、七年）其间往来文书，证明「浩罕自来不知尊卑，习俗如此」，认为「其宗教叔伯尊卑尚自不失所知，不值与之计较」27，才使浩罕使者得以进京朝觐。

由以上史料所示，可知从十八世纪中期，浩罕显然成了清朝的藩属国。而且其关系之亲密，就如同霍渥斯所说:「纳尔巴图平均每一年或二年向清朝朝贡一次；乾隆皇帝也称他为自己的儿子。」28

基本上，自额尔德尼以迄爱里木三代伯克间，清政府对其一方面在政治上只要「一切事物俱遵驻扎喀什噶尔、叶尔羌大臣等节制」29即可，而另一方面使节朝觐时也款待礼遇有嘉，「今既为臣仆，自必加恩抚恤」30多所赏赐，「加恩召见，频申宴赉，赐尔缎匹，器什若干。尔其祇受，又赏来使及众人等服饰、器用、银两有

\vfill
\begin{BookNotes}
\BookNote{25.}{傅恒，《钦定皇舆西域图志》(台北，文海出版社，民国59年)，页13。}
\BookNote{26.}{《清高宗实录》卷601，页7，乾隆二十四年十一月甲子谕。}
\BookNote{27.}{章佳容安，《那文毅公奏议》卷19（台北，文海出版社，民国55年），页418。}
\BookNote{28.}{Henry Hoyle Howarth, History of the Mongols, 4vols, (London : Longmans, Green, \&Co., 1876-1927) , Vol. 3, p.818.}
\BookNote{29.}{《清高宗实录》卷678，页17，乾隆二十八年正月己巳条。}
\BookNote{30.}{同上注，卷606，页2，乾隆二十五年二月丙子条。}
\end{BookNotes}

% 原书第 265 页（PDF 第 278 页）
\BookPage{265}{278}
\noindent 差」31。浩罕来使，只要政治上之臣服，以象征性之贡物入京朝觐，即可得到清政府颇丰的回赏，甚至有准商减税的优待，「外来贸易之人，征二十分之一，每岁俱未能如数交纳，臣等酌量减赋」32。当然入觐之同时，其使节也沿途买进大量货物。因此在这种获得清朝政府的统治「朝贡」与贸易互通有无的有利条件下，浩罕自爱里木称汗至艾玛尔（Omar，1811-1822）时，汗国的国家机构建立，汗国的刍形已成。根据伊凡诺夫的《中亚史纲》中说:「在爱尔玛时代，官吏总数达四万人。」33纳利夫金也说：「经过一翻组织以后的浩罕汗国具备了国家机构的完整形式。」34

\BookSubheading{四、清朝中衰，居奇「逆裔」因以坐大}

十九世纪初，清代中国本部，因为政治和军队的腐败，人口急遽增加之压力与汉人排满的风潮云涌，而发生一连串的「会匪乱党」事件，正显现出清朝统治势力全面发生危机的征兆。在边疆地区的局势，也因此而逐渐受着这种国势衰颓的影响，诸多的冲突与矛盾，也渐渐的表露出来。当时新疆的天山北路方面，因为主要是属蒙古族牧地，民族比较单纯，而且清朝在平准噶尔时，准部伤亡惨重，「千里空虚，渺无人烟」35，至于明代西迁俄境额济勒河，后因不胜俄对其征兵之苦的土尔扈特部，于乾隆三十六年（1771）返抵伊犁「卡伦」36时，仅存七万余口，因分散北疆各地，人数不多，所以「乱源肃清，乱事即无从发生」37。只有南路回疆地区，由于清朝在新疆驻军，一向重北轻南，且该地区民族也较复杂，又有所谓大小和卓「逆裔」外逃，时

\vfill
\begin{BookNotes}
\BookNote{31.}{同上注，页1。}
\BookNote{32.}{同上注，卷593，页11，乾隆二十四年七月庚午条。}
\BookNote{33.}{伊凡诺夫，前揭书，页85，引自潘志平，前揭书，页67。}
\BookNote{34.}{纳利夫金，《浩罕国简史》（巴黎，1889年)，引自潘志平，前揭书，页67。}
\BookNote{35.}{椿园七十一，《西域总志》卷3（台北，文海出版社，民国55年)，页5。}
\BookNote{36.}{卡伦（Karun)一词，源于古突厥语，出现在我国历史上的辽代，后为女真族沿用，直至清代译为汉语「卡伦」，满洲语，为斥候之意。而突回语Kharam，Qaram，最古的原意是武器的描准星，后来演变为瞭望，哨或哨所。何秋涛《朔方备乘》卷10〈北徼卡伦考〉中说：「更番候望之所曰台，国语谓之喀伦。」其中所谓「国语」乃指清代之满洲语而言。详见于福顺，〈清代新疆卡伦述略〉《历史研究》第四期（北京，中国社会出版社，1979)，页75；及刘义棠，中国边疆民族史上课讲义。}
\BookNote{37.}{曾问吾，《中国经营西域史》，中册（台北，文海出版社，民国25年影印)，页298。}
\end{BookNotes}

% 原书第 266 页（PDF 第 279 页）
\BookPage{266}{279}
\noindent 时有恢复其「祖父地业」38的复辟企图，且回众笃信经义，对于宗教极为狂热，或由个人野心，加上偶有外人助力，即行变作。而此时浩罕正自汗国刍形骤成后，即趁清朝之中衰，边务守备的放松，如「以节省兵力为由，命换防部队兼任巡边任务，至使巡边越走越近，甚至放弃某些地方的巡防，且移设和添设卡伦也逐渐废止」39，而逐渐「叛离」清朝。当爱玛尔继立其兄的事业（浩罕汗谱系如附表一），向东扩张，占据鄂什，该地原为清朝所属布鲁特额德格纳部牧地。如今额德格纳部不得不退居噶布兰、苏提布拉克，并以此二岭与之为界，就如徐松《西域水道记》卷一所说的「岭东为回部，领西为霍罕」了，所以至此浩罕势力已达清代中国回疆边外。

浩罕自爱里木称汗，国势日强。大和卓之子萨木萨克（Samsaq）匿遁其境，因此居为奇货，如清史稿所载，大和卓「波罗尼都次子萨木萨克逃入浩罕，藉其和卓之名，居为奇货」40，并于嘉庆十四年（1890），以攻下塔什干，向清廷报喜为由，借此夸耀且要挟入卡免税，「每次烦碎，恳求免税」41。直至道光二年（1822），爱玛尔子迈买底里（MuhammadAli，1822-1842）即位，即开始以喀什噶尔和卓后裔，在喀什噶尔发动有组织的「圣战」42，「作为教胞的浩罕汗国，也带着政治的、经济的意图，予以支援」43。浩罕从此因而坐大，正式与清政府决裂，并多次利用回疆人民对中亚传统宗教情感之强烈认同与败逃于其境内之和卓后裔，进行勾结，开始对中国回疆入卡侵扰。

\vfill
\begin{BookNotes}
\BookNote{38.}{曹振镛，《平定回疆剿擒逆裔方略》卷63(台北，文海出版社，民国55年)，页26。}
\BookNote{39.}{杨建新，〈清代的西北边防政策和沙俄侵略〉《西北史地》第一期（兰州，兰州大学出版社，1980年），页37。}
\BookNote{40.}{赵尔巽，《清史稿》卷536，属国4，浩罕（台北，国史馆校注本，民国79年），页12155。}
\BookNote{41.}{章佳容安，前揭书，卷19，页14。}
\BookNote{42.}{「圣战」（Djihad)：阿拉伯语，原意为「奋斗」，即谓「为安拉之道」而战，详见郑勉之，《伊斯兰教简明辞典》（南京，江苏古籍出版社，1993年），页260。又圣战，即吉哈德（djihad）源于阿拉伯文动词djahada，为奋斗尽力之意，其宗教含意，有为传播伊斯兰信仰而奋斗的意义。在伊斯兰教义中，圣战是一种义务，且是真主指引到达天堂最直接之路。详见M. Khaddus,Was andpeacein the law of Islam(TheJohnsHopkins:1955)，pp.56-57及萧宪，〈伊斯兰教吉哈德理论在当代的发展和演变〉，《世界宗教研究》二期（北京，中国社会科学出版社，1988年6月），页32。}
\BookNote{43.}{佐口透着，凌颂纯译，《十八一十九世纪新疆社会史研究》下册（乌鲁木齐，新疆人民出版社，1983年），页504。}
\end{BookNotes}

% 原书第 267 页（PDF 第 280 页）
\BookPage{267}{280}
五结论

汉武帝通西域，纳天山南北为中国之城垣，处其毗邻中亚为中国之外郭；清高宗荡平准部、大小和卓，统一天山南北，赐名为新疆，使得原本只为得其土不足以耕耘，得其民不足以驱使的古西胡蛮荒之域，顿时成为屏障中国西北，兼司北徼的国家根本要地。无奈浩罕凭其所处的自然天险，居欧亚陆路交通枢纽之费尔于纳盆地的地理优越条件，并在中亚布哈拉、希瓦为争河中地区统治权而不断争战的狭缝中乘机脱离其管辖。且利用准部被清朝所平定，使长期以来，一面忍受其侵凌，并一面巩固风雨飘摇中所建立起来的政权，骤解其桎梏。同时借助附庸清代中国通贡受赐的机会而崛起，最后终于在十九世纪的二十年代里，趁着清朝政府之中衰，居奇和卓后裔，要挟清政府通商免税而壮大。从此使得回疆无一宁日，甚至关陇险些隘其封，河湟也渐失其险。直迄今日，它当时的崛起与侵扰，除所对近代中国回疆的治理措施，无论是在宗教层面或是民族关系的调整，影响深远外，吾人更应在产生过的相当震撼的历史变作中记取教训，对外交涉策略上也委实足资让我们对之深省与借鉴。


% 原书第 268 页（PDF 第 281 页）
\BookPage{268}{281}
\BookArt{p281-01.jpg}{0.970}{0.829}


% 原图说　附表一：浩罕汗谱系表
% 原图说　录自潘志平，《中亚浩罕国与清代新疆》（北京，中国社会科学出版社，1991年），页177。
